\documentclass{amsart}
% For dealing with references we use the comment environment
\usepackage{verbatim}
\newenvironment{reference}{\comment}{\endcomment}
%\newenvironment{reference}{}{}
\newenvironment{slogan}{\comment}{\endcomment}
\newenvironment{history}{\comment}{\endcomment}

% For commutative diagrams we use Xy-pic
\usepackage[all]{xy}

% We use 2cell for 2-commutative diagrams.
\xyoption{2cell}
\UseAllTwocells

% We use multicol for the list of chapters between chapters
\usepackage{multicol}

% This is generally recommended for better output
\usepackage{lmodern}
\usepackage[T1]{fontenc}

% For cross-file-references

% Package for hypertext links:
\usepackage{hyperref}

% For any local file, say "hello.tex" you want to link to please

% Theorem environments.
%
\theoremstyle{plain}
\newtheorem{theorem}[subsection]{Theorem}
\newtheorem{proposition}[subsection]{Proposition}
\newtheorem{lemma}[subsection]{Lemma}

\theoremstyle{definition}
\newtheorem{definition}[subsection]{Definition}
\newtheorem{example}[subsection]{Example}
\newtheorem{exercise}[subsection]{Exercise}
\newtheorem{situation}[subsection]{Situation}

\theoremstyle{remark}
\newtheorem{remark}[subsection]{Remark}
\newtheorem{remarks}[subsection]{Remarks}

\numberwithin{equation}{subsection}

% Macros
%
\def\lim{\mathop{\mathrm{lim}}\nolimits}
\def\colim{\mathop{\mathrm{colim}}\nolimits}
\def\Spec{\mathop{\mathrm{Spec}}}
\def\Hom{\mathop{\mathrm{Hom}}\nolimits}
\def\Ext{\mathop{\mathrm{Ext}}\nolimits}
\def\SheafHom{\mathop{\mathcal{H}\!\mathit{om}}\nolimits}
\def\SheafExt{\mathop{\mathcal{E}\!\mathit{xt}}\nolimits}
\def\Sch{\mathit{Sch}}
\def\Mor{\mathop{\mathrm{Mor}}\nolimits}
\def\Ob{\mathop{\mathrm{Ob}}\nolimits}
\def\Sh{\mathop{\mathit{Sh}}\nolimits}
\def\NL{\mathop{N\!L}\nolimits}
\def\CH{\mathop{\mathrm{CH}}\nolimits}
\def\proetale{{pro\text{-}\acute{e}tale}}
\def\etale{{\acute{e}tale}}
\def\QCoh{\mathit{QCoh}}
\def\Ker{\mathop{\mathrm{Ker}}}
\def\Im{\mathop{\mathrm{Im}}}
\def\Coker{\mathop{\mathrm{Coker}}}
\def\Coim{\mathop{\mathrm{Coim}}}

% Boxtimes
%
\DeclareMathSymbol{\boxtimes}{\mathbin}{AMSa}{"02}

%
% Macros for moduli stacks/spaces
%
\def\QCohstack{\mathcal{QC}\!\mathit{oh}}
\def\Cohstack{\mathcal{C}\!\mathit{oh}}
\def\Spacesstack{\mathcal{S}\!\mathit{paces}}
\def\Quotfunctor{\mathrm{Quot}}
\def\Hilbfunctor{\mathrm{Hilb}}
\def\Curvesstack{\mathcal{C}\!\mathit{urves}}
\def\Polarizedstack{\mathcal{P}\!\mathit{olarized}}
\def\Complexesstack{\mathcal{C}\!\mathit{omplexes}}
% \Pic is the operator that assigns to X its picard group, usage \Pic(X)
% \Picardstack_{X/B} denotes the Picard stack of X over B
% \Picardfunctor_{X/B} denotes the Picard functor of X over B
\def\Pic{\mathop{\mathrm{Pic}}\nolimits}
\def\Picardstack{\mathcal{P}\!\mathit{ic}}
\def\Picardfunctor{\mathrm{Pic}}
\def\Deformationcategory{\mathcal{D}\!\mathit{ef}}

\setlength{\emergencystretch}{3em}
\begin{document}
\title[Stacks Project, Tag 08GP]{390 --- Nisnevich induction principle for quasi-compact quasi-separated algebraic spaces}
\maketitle
\noindent Copyright (C) 2005--2025 Johan de Jong. Stacks Project authors. Source: \href{https://stacks.math.columbia.edu/tag/08GP}{Tag 08GP}.

\noindent Editorial difficulty note (not author text). Difficulty H2: An algebraic space cannot be covered by affine Zariski opens, so the scheme induction principle does not suffice. The proof uses a finite scheme-stratified etale filtration and refines each stage into elementary distinguished squares with affine chart, showing that all overlap objects meet the inductive hypotheses..

\noindent Editorial provenance note (not author text). History: excerpt from revision \texttt{a0444\allowbreak 6e57e\allowbreak c1fbc\allowbreak 252a8\allowbreak 71afc\allowbreak ec775\allowbreak 2fb28\allowbreak 07b14}, spaces-perfect.tex, lines 1345--1423. Reference presentation was adapted; original-author context and core support blocks were retained or added. The mathematical wording is unchanged. Licensed under GNU FDL 1.2 or later; no invariant sections or cover texts. See ../licenses/COPYING. Context blocks reproduce author definitions and supporting results; they are not additional counted problems.

\begin{definition}
\label{definition-elementary-distinguished-square}
Let $S$ be a scheme. A commutative diagram
$$
\xymatrix{
U \times_W V \ar[r] \ar[d] & V \ar[d]^f \\
U \ar[r]^j & W
}
$$
of algebraic spaces over $S$ is called an {\it elementary distinguished square}
if
\begin{enumerate}
\item $U$ is an open subspace of $W$ and $j$ is the inclusion morphism,
\item $f$ is \'etale, and
\item setting $T = W \setminus U$ (with reduced induced
subspace structure) the morphism $f^{-1}(T) \to T$ is an isomorphism.
\end{enumerate}
We will indicate this by saying: ``Let $(U \subset W, f : V \to W)$
be an elementary distinguished square.''
\end{definition}

\begin{lemma}
\label{lemma-induction-principle}
Let $S$ be a scheme. Let $X$ be a quasi-compact and quasi-separated
algebraic space over $S$. Let $P$ be a property of the quasi-compact
and quasi-separated objects of $X_{spaces, \etale}$. Assume that
\begin{enumerate}
\item $P$ holds for every affine object of $X_{spaces, \etale}$,
\item for every elementary distinguished square $(U \subset W, f : V \to W)$
such that
\begin{enumerate}
\item $W$ is a quasi-compact and quasi-separated object of
$X_{spaces, \etale}$,
\item $U$ is quasi-compact,
\item $V$ is affine, and
\item $P$ holds for $U$, $V$, and $U \times_W V$,
\end{enumerate}
then $P$ holds for $W$.
\end{enumerate}
Then $P$ holds for every quasi-compact and quasi-separated object
of $X_{spaces, \etale}$ and in particular for $X$.
\end{lemma}

\begin{proof}
We first claim that $P$ holds for every representable
quasi-compact and quasi-separated object of $X_{spaces, \etale}$.
Namely, suppose that $U \to X$ is \'etale and $U$ is a
quasi-compact and quasi-separated scheme. By assumption (1)
property $P$ holds for every affine open of $U$. Moreover, if
$W, V \subset U$ are quasi-compact open with $V$ affine and $P$ holds for 
$W$, $V$, and $W \cap V$, then $P$ holds for $W \cup V$ by (2)
(as the pair $(W \subset W \cup V, V \to W \cup V)$ is an elementary
distinguished square). Thus $P$ holds for $U$ by the induction
principle for schemes, see
Cohomology of Schemes, Lemma \href{https://stacks.math.columbia.edu/tag/08DR}{lemma induction principle (Tag 08DR)}.

\medskip\noindent
To finish the proof it suffices to prove $P$ holds for $X$
(because we can simply replace $X$ by any quasi-compact and quasi-separated
object of $X_{spaces, \etale}$ we want to prove the result for).
We will use the filtration
$$
\emptyset = U_{n + 1} \subset
U_n \subset U_{n - 1} \subset \ldots \subset U_1 = X
$$
and the morphisms $f_p : V_p \to U_p$ of
Decent Spaces, Lemma
\href{https://stacks.math.columbia.edu/tag/07ST}{lemma filter quasi compact quasi separated (Tag 07ST)}.
We will prove that $P$ holds for $U_p$ by descending induction on $p$.
Note that $P$ holds for $U_{n + 1}$ by (1)
as an empty algebraic space is affine. Assume $P$ holds for $U_{p + 1}$.
Note that $(U_{p + 1} \subset U_p, f_p : V_p \to U_p)$ is an elementary
distinguished square, but (2) may not apply as $V_p$ may not be affine.
However, as $V_p$ is a quasi-compact scheme we may choose a finite affine open
covering $V_p = V_{p, 1} \cup \ldots \cup V_{p, m}$.
Set $W_{p, 0} = U_{p + 1}$ and
$$
W_{p, i} = U_{p + 1} \cup f_p(V_{p, 1} \cup \ldots \cup V_{p, i})
$$
for $i = 1, \ldots, m$. These are quasi-compact open subspaces of $X$.
Then we have
$$
U_{p + 1} = W_{p, 0} \subset
W_{p, 1} \subset \ldots \subset
W_{p, m} = U_p
$$
and the pairs
$$
(W_{p, 0} \subset W_{p, 1}, f_p|_{V_{p, 1}}),
(W_{p, 1} \subset W_{p, 2}, f_p|_{V_{p, 2}}),\ldots,
(W_{p, m - 1} \subset W_{p, m}, f_p|_{V_{p, m}})
$$
are elementary distinguished squares by
Lemma \href{https://stacks.math.columbia.edu/tag/08GN}{lemma make more elementary distinguished squares (Tag 08GN)}.
Note that $P$ holds for each $V_{p, 1}$ (as affine schemes) and for
$W_{p, i} \times_{W_{p, i + 1}} V_{p, i + 1}$ as this is a quasi-compact
open of $V_{p, i + 1}$ and hence $P$ holds for it by the first paragraph
of this proof. Thus (2) applies to each of these and we inductively
conclude $P$ holds for $W_{p, 1}, \ldots, W_{p, m} = U_p$.
\end{proof}
\end{document}
